\chapter{认知相干动力学：绑定、粗粒化与无中心协同}
\label{chp:cognitive_coherence_dynamics}

上一章把运行中的边流分解为任务推进、循环维持和全局约束，本章进一步讨论多个局部状态怎样形成可供任务使用的整体表示。我们把“相干”定义为一组带身份的局部记录，在给定查询、时间窗和干预集合下，能够维持一致绑定并产生稳定读出的性质；它不是所有单元振幅相同，也不是系统进入某种神秘的全域共振。周期变量之间的同步可以成为相干证据，却既非必要条件，也非充分条件。为此，我们先建立局部耦合与对象绑定模型，再说明任务相关序参量如何从粗粒化中产生，随后把全局循环与结构冗余纳入保持条件，最后分析没有宏观控制器时的自组织协同及其失效边界。整章始终区分定义、模型假设、数学结果和工程验收，并把局部结构网络与全局分布表示放在同一套身份、绑定、耦合和读出接口中考察。

\section{局部耦合：同步证据、对象身份与绑定提交}
\label{sec:local_dynamics_phase_locking}

设运行时包含$n$个局部单元，单元$i$的记录为$x_i=(v_i,a_i,\iota_i,s_i,\nu_i)$。其中$v_i\in\mathbb R^{d_i}$是内容或特征状态，$a_i\geq0$是无量纲活动强度，$\iota_i$是对象或事件身份键，$s_i$记录来源与时间，$\nu_i$是版本。局部结构网络$G_t=(V_t,E_t)$保存角色边、兼容约束和候选绑定；全局分布表示$p_t(z)$表达对象类别、解释或行动候选的不确定性。两者不能靠“相位相同”自动合并：红色区域与圆形轮廓可能同步响应同一刺激，却属于相邻的两个对象；同一对象的颜色和位置也可能由不同采样率的通道给出，从不形成稳定相位差。我们因此把绑定候选写成$b=(i,j,r,\iota,\rho)$，其中$r$是角色关系，$\iota$是共同身份假设，$\rho$是来源可追溯记录。只有身份兼容、时空支持、角色约束和反事实检验共同通过时，候选才能提交到结构网络。

在某些实现中，局部单元确实具有周期或准周期状态，例如扫描控制器、振荡采样器、轮询工作集或分时通信协议。此时可令$\theta_i\in\mathbb S^1$表示相位，$\omega_i$表示每单位墙钟时间的本征更新率，$K_{ij}\geq0$表示从$j$到$i$的无量纲耦合增益，$\tau_{ij}\geq0$表示传输时延，$\varphi_{ij}$表示接口协议要求的相位偏置。一个可检验的延迟耦合模型是
\begin{equation}
\dot\theta_i(t)=\omega_i(t)+\sum_{j\in\mathcal N_i(t)}K_{ij}(t)a_j(t)
\sin\!\left[\theta_j(t-\tau_{ij})-\theta_i(t)-\varphi_{ij}(t)\right]+\xi_i(t),
\label{eq:task_conditioned_phase_coupling}
\end{equation}
其中$\xi_i$是零均值扰动的模型项，而不是“认知热量”。这个方程是模型假设，不是从目的论场方程推导出的普遍规律。若单元没有可重复周期、相位估计受混叠影响，或事件驱动过程只有离散时间戳，就应改用消息到达、因果顺序和版本一致性模型。即使相位模型适用，$K_{ij}$也应由通信可靠性、任务相关性和许可策略估计，不能把概念嵌入距离直接称为物理导通率；$\varphi_{ij}$是接口或时序偏置，不是由价值规范场产生的贝里相位。

在固定图、固定参数、无时延且对称耦合的专门情形，令$e_{ij}=\theta_j-\theta_i-\varphi_{ij}$。若存在一组一致的节点偏置使所有$\varphi_{ij}$可由节点势差表示，则局部锁定状态满足$\dot e_{ij}=0$，其稳定性可由线性化后的加权图拉普拉斯判断。这个数学结果只说明相对时序可能稳定；它不证明两个内容属于同一对象，更不证明内容真实。反过来，绑定也不要求锁相：异步视觉区域、语言指称和外部数据库记录可以通过共同身份键、时间容差和来源链完成绑定。相干判定因此采用证据向量
\begin{equation}
c_{ij}=(c_{\mathrm{id}},c_{\mathrm{time}},c_{\mathrm{role}},c_{\mathrm{source}},c_{\mathrm{intervene}},c_{\mathrm{phase}}),
\end{equation}
先以身份、权限和冲突规则筛除不合法候选，再由任务特定读出器综合其余证据。相位一致性只是最后一个可选分量。

以“红色圆球”为例，系统分别得到颜色区域、轮廓、位置和语言指称。若两个相邻球同时被红光照亮，颜色与轮廓的共同激活会制造虚假同步。正确流程应追踪分割掩码和对象身份，遮蔽一个球或移动摄像机后重新测量：若颜色证据随同一身份迁移且竞争对象的解释概率下降，绑定才获得方向性支持。提交时还应保留未选候选和置信区间，以便新观测到达后回滚。AgentOS可以让$h(t)$保存当前候选，$E$保存带来源的绑定事件，$\Theta$保存长期角色约束，SPC形成任务相关压缩，TCE调节候选竞争，Boundary限制跨源写入，外部记忆$M_e$保留原始证据；这些只是工程实例。一般HSF-HD模型只要求局部记录、身份接口、耦合证据、可撤销提交和外部验收，并不规定所有智能系统必须采用这些组件。

\section{全域粗粒化：任务序参量与快慢变量约化}
\label{sec:global_dynamics_order_parameter}

当局部单元数量很大时，逐项检查全部状态既昂贵也不利于控制。我们定义任务$q$、观测窗口$W$和候选集合$C_q$，再由粗粒化算子$\mathcal C_q$生成宏观状态$z_q=\mathcal C_q(G_t,\{x_i\},p_t)$。所谓序参量不是脱离任务客观存在的“全局意义”，而是对某组查询和预测足够的低维统计量。若单元具备可信相位，可定义带权相干量
\begin{equation}
Z_q(t)=R_q(t)e^{\mathrm i\Theta_q(t)}=
\frac{1}{\sum_{i\in C_q}w_i(t)}\sum_{i\in C_q}w_i(t)e^{\mathrm i(\theta_i(t)-\chi_i(t))},
\label{eq:task_order_parameter}
\end{equation}
其中$w_i\geq0$综合来源可靠性、任务相关性和身份资格，$\chi_i$是已声明的接口对齐偏置。$R_q\in[0,1]$只度量这个候选集合在所选相位坐标中的集中程度；集合、权重或偏置改变时，$R_q$的变化不能全部解释为系统动力学。没有相位结构的系统可使用一致投票率、后验熵、查询失真、约束满足率或消息版本收敛度作为序参量，不能为了套用复数公式而虚构振子。

宏观量能否支配局部变量，需要额外的快慢尺度条件。令$u\in\mathbb R^m$表示快速局部偏差，$z\in\mathbb R^r$表示慢速宏观变量，并考虑
\begin{equation}
\epsilon\dot u=f(u,z,q,t),\qquad \dot z=g(u,z,q,t),\qquad 0<\epsilon\ll1.
\label{eq:coherence_fast_slow_system}
\end{equation}
若对每个允许的$z$，方程$f(u,z,q,t)=0$存在局部唯一解$u=h(z,q,t)$，且快速子系统在相关区域内一致指数稳定，那么在有限时间窗内可用$u\approx h(z,q,t)$约化原系统。协同论中的“伺服”在这里是一条带条件的模型约化结论：局部自由度没有被消灭，只是在误差界内迅速靠近慢流形。若结构切换、时延增长、输入突变或快速子系统出现多个吸引支，唯一慢流形假设就会失败；系统必须保留局部状态和模式编号，不能让宏观读出覆盖冲突证据。

为了描述宏观状态的形成，可在固定结构区间采用正规形
\begin{equation}
\tau_z\dot z=\alpha(q,u_{\mathrm{ext}},t)z-\beta\|z\|^2z-DL_qz+b_q(t),
\label{eq:coherence_normal_form}
\end{equation}
其中$L_q$是任务相关图拉普拉斯，$D\geq0$控制邻域平滑，$b_q$是外部证据，$\alpha$和$\beta>0$是经验辨识参数。固定$\alpha>0$且忽略网络与输入时，非零平衡模长为$\sqrt{\alpha/\beta}$；这只是正规形的数学结果，不能据此把心流、顿悟或理解宣布为相变。真实系统中目标$q$、结构$L_q$、输入$b_q$和度量都可能变化，因此评价函数$V(z,q,L,b,t)$的总导数必须包含$\partial_qV\dot q$、$\langle\partial_LV,\dot L\rangle$、$\langle\partial_bV,\dot b\rangle$和$\partial_tV$。局部收敛、预测误差下降、格式塔稳定和任务成功是四个不同结论，需要分别验收。

宏观相干也可能过强。多个单元迅速同意，可能来自共享错误输入、复制污染、中心节点压制或候选多样性丢失。我们因此同时记录$R_q$、有效独立来源数、身份冲突率、校准误差、少数候选覆盖率和干预敏感度。对多Agent协作，消息一致不等于事实正确；对语言理解，词元激活集中不等于指称已经确定；对控制系统，命令同步不等于执行器处于安全状态。验证时可遮蔽共同输入、移除高中心度节点、交换身份键、延迟一个通道并引入反事实样本。只有宏观读出在允许扰动下保持任务性能，同时能够暴露不一致和拒绝错误提交，才能称为有效相干。局部结构网络提供身份与约束，全局分布表示提供不确定性和候选覆盖，序参量只是两者之上的任务读出，不取代任何一方。

粗粒化还必须处理宏观读出的滞后与回写。若$z_q$由长度为$W$的滑动窗口估计，它天然落后于最新局部事件；控制器若据此提高耦合，可能在任务已经切换后继续强化旧簇。实现时应为每个序参量记录估计窗口、有效样本量、更新时间和适用任务版本，并把回写操作限制在同一事务版本。对于多个竞争宏观解释，系统不应过早平均为单一$z_q$，而应保持带权候选$(z_q^{(1)},\ldots,z_q^{(m)})$，直到区分性观测到达。验收不仅比较压缩前后的预测性能，还应测量读出延迟、回写振荡、模式遗漏和从错误宏观状态恢复所需的时间。这些指标使“整体格式塔”成为可复现的工程对象，而不是由高相干度自动赋予的现象学名称。

\section{全局约束：循环空间、结构冗余与意义保持}
\label{sec:topological_endpoint_harmonic_flow}

上一章已经说明，离散Hodge分解中的调和分量依赖所选复形、权重和边界条件。本章关心它怎样参与相干保持，而不是把调和零模解释为零耗散或“意义结晶”。设任务相关胞腔复形为$\mathcal K_q$，边流$y\in C_1(\mathcal K_q)$记录证据、消息或控制量的有向传递，边界算子为$B_1$和$B_2$。在加权内积下，满足$B_1y=0$且$B_2^\ast y=0$的分量属于调和子空间。它表示不能由局部节点势差或面内循环消去的全局循环自由度。这样的分量可能对应跨模块承诺、长期依赖、未填补的结构孔洞，也可能只是数据缺失和建模分辨率造成的伪循环。代数上的核空间并不产生任务含义，只有当其与身份、来源、查询和干预结果建立稳定对应后，它才成为可用的全局约束。

我们把“意义保持”操作化为查询保持。给定查询族$\mathcal Q$、允许扰动集合$\mathcal P$和读出器$r_q$，若压缩或局部损坏后的状态$X'$满足
\begin{equation}
\sup_{q\in\mathcal Q,\,p\in\mathcal P}
d_q\!\left(r_q(p(X')),r_q(p(X))\right)\leq\varepsilon_q,
\label{eq:meaning_query_preservation}
\end{equation}
则称表示在该查询族上以误差$\varepsilon_q$保持。这里$d_q$是任务声明的输出距离；分类、轨迹、文本和安全约束需要不同度量。这个定义允许意义分布在多个局部记录和全局参数中，却不宣称其编码在不可破坏的拓扑干涉条纹里。删除部分节点后能够恢复，可能来自纠错码、冗余证据、生成先验或外部记忆，而不必来自同调不变量；反之，图的贝蒂数不变也不能保证对象身份、数值属性或规范约束得到保留。

结构冗余需要区分有益备份与相关失败。若三个副本都来自同一个错误传感器，它们增加数量却不增加独立证据；若多个路径共享同一个权限服务，它们在拓扑上不同却会同时失效。系统应为每条恢复路径记录来源祖先、故障域、更新时间和适用查询。我们可构造恢复矩阵$M_{qk}$，表示组件$k$失效时查询$q$能否在误差阈值内恢复，并以最坏列损失、相关故障模拟和恢复时延评估鲁棒性。调和分量较大时，优先检查是否存在没有局部责任节点的全局承诺或环形依赖；调和分量消失时，也不能断言系统没有全局意义，因为树形结构和集中式数据库同样能稳定回答查询。拓扑诊断提供结构线索，不替代语义和因果验证。

“无旋、无散”同样不等于无计算损耗。即使边流位于离散调和子空间，维持消息、校验副本、刷新缓存和运行硬件仍消耗墙钟时间、内存带宽与焦耳能量。反过来，一个包含局部散度的系统可能正从传感器吸收新证据，包含局部旋度的系统可能在执行必要复核。我们分别建立活动账本$A$、任务代价$J$、资源用量$C$和设备能耗$E_{\mathrm{J}}$，只通过实测模型连接它们，不以同一个“认知能量”替代。工程验收包含查询保持、故障恢复、证据独立性、版本一致性和资源曲线；任何一项都不能由$\Delta y=0$单独推出。

在AgentOS实例中，$E$可保存可追溯事件环，$\Theta$保存跨时间约束，$h(t)$承载当前读出，CCM检测循环占用和结构缺口，Offloading把冗余证据迁移到$M_e$，Boundary维护写入权限，SPC和TCE选择当前任务的压缩与控制强度。若外部记忆不可达，系统应降级到明确的保留查询集，而不是凭相干感继续回答。一般智能系统也可用数据库事务、因子图、纠错编码、黑板架构或联邦协议实现同类功能。HSF-HD关心的是全局约束怎样限制状态轨迹、局部损坏怎样传播以及读出怎样被验证，不把某种拓扑或记忆架构规定为唯一实现。

\section{无中心协同：自组织阈值、伪相干与调控边界}
\label{sec:resonance_without_macro_layer}

宏观控制器不是形成协同的必要条件。固定无向连通图上，若$K_{ij}=K_{ji}\geq0$、无时延、相位偏置为零且本征频率相同，耦合动力学可写成势函数
\begin{equation}
U(\theta)=-\frac12\sum_{i,j}K_{ij}\cos(\theta_i-\theta_j),
\qquad \dot\theta_i=-\frac{\partial U}{\partial\theta_i}.
\label{eq:decentralized_coherence_potential}
\end{equation}
于是$dU/dt=-\sum_i(\partial U/\partial\theta_i)^2\leq0$，系统在给定吸引域内趋向势函数的临界集合。这个数学结果说明局部规则可能产生群体秩序，但它依赖对称静态耦合、共同频率和连续更新。异质频率、有向非正规图、通信时延、符号相反的耦合、结构切换和噪声都会改变阈值与稳定性；存在多个局部极小值时，系统还可能收敛到分簇、扭曲或错误一致状态。由局部下降不能推出全局同步，更不能推出事实、价值或任务正确。

在随机振子总体的经典均场近似中，临界耦合与频率分布在中心附近的密度有关；在有限稀疏网络中，谱间隙、度异质性和噪声会共同影响相干形成。我们不把单一$K_c$当作所有认知系统的普遍常数，而将阈值视为需要实验辨识的函数$K_c(G,g_\omega,\tau,\sigma,q)$。测试应扫描耦合、噪声和时延，报告迟滞、初值敏感性和有限规模置信区间。对于非振荡系统，对应问题是局部消息规则何时使信念、版本或承诺收敛，可使用一致性算法和异步分布式优化分析；其成立条件涉及图的联合连通、消息可靠性和目标兼容性，而不是“平均场化身为幽灵中心”。宏观序参量是观察者或控制器计算出的统计量，不是新的主体，也没有脱离局部接口的绝对统治力。

无中心协同既能产生有效秩序，也能放大错误。耳虫式重复可建模为高自回返增益和低新信息输入造成的局部吸引；自发联想链可能提供候选，也可能占用工作集而不增加任务进展；群体转发可能形成稳定共识，也可能源于共同谣言。诊断不能只看$R_q$上升，而要联合新信息率、来源独立性、任务改善、退出能力和反事实响应。若遮蔽一个共同输入后相干立即消失，则它主要由外源驱动；若打断少数高中心边即可恢复多样性，则问题在耦合拓扑；若所有单元在不同证据下仍保持同一结论，则可能存在硬编码偏置。软件指标只能诊断计算过程，不能据此给人类作临床判断。

分布式实现还要防止拜占庭节点、陈旧版本和网络分区。只有普通崩溃故障时，多数确认可能足以推进；存在恶意或系统性相关错误时，简单多数和相位一致都不构成可靠证明。节点应签名消息、携带任务与版本标识，并把超时、不确定和拒绝列为合法输出。网络重新连通后，需要按照因果顺序合并承诺，对冲突执行显式仲裁，不能把最后到达或幅值最大的状态视为自然真值。由此可见，无中心不等于无治理：治理可以分布在协议、权限、证据门槛和回滚规则中，而不必集中于单一宏观主体。

调控的目标不是最大化同步，而是在任务需要的时间和范围内形成足够相干，同时保留异议、探索和退出路径。控制器可降低错误回返边的增益、提高新证据权重、拆分竞争簇、延长观测窗口、请求外部验证或暂缓提交。变化量若通过离散实现施加，应声明采样步长、通信周期、饱和上限和回滚版本；连续模型的稳定区间不能直接移植到异步部署。对于不可逆外部行动，相干只提供提交证据之一，还必须通过权限、安全和现实可行性检查。由此，我们得到本章的限定性结论：相干是任务相关的联合可用性，可能借助同步但不等同同步；宏观变量可以压缩和调控局部状态，但不自动代表整体真理；无中心系统能够自组织，也会形成伪相干。只有身份绑定、结构约束、分布不确定性、干预测试和外部验收共同闭环，局部活动才真正成为可执行的整体认知状态。

%--chp--
